\chapter{Conclusioni e Sviluppi Futuri}
\label{ch:conclusioni}

\section{Conclusioni}

L'obiettivo di questo lavoro era la progettazione e l'implementazione di una piattaforma IoT integrata per il trasporto refrigerato condiviso, capace di unire in un unico sistema due funzionalità tradizionalmente separate: un marketplace per il matching tra domanda e offerta di spedizioni e un sistema di monitoraggio delle condizioni del vano di refrigerazione dei veicoli con rilevazione automatica delle anomalie.


\section{Sviluppi futuri}

Nonostante i risultati ottenuti, il lavoro presenta margini di miglioramento che possono essere affrontati in sviluppi successivi.

\subsection{Ottimizzazione temporale e scheduling attivo}

L'attuale algoritmo di matching tra spedizioni e viaggi opera prevalentemente sulla dimensione spaziale, valutando la prossimità geografica tra i punti di ritiro/consegna e la polyline del percorso. La dimensione temporale, pur presente sotto forma di data di scadenza della spedizione, non viene sfruttata in modo strutturato per guidare la pianificazione dei viaggi. Un sistema più maturo potrebbe integrare le finestre di disponibilità dei trasportatori, i periodi di inattività dei veicoli e i vincoli temporali delle spedizioni in un unico modello di ottimizzazione.
\subsection{Scalabilità su piattaforme cloud}

L'attuale architettura è validata in ambiente locale tramite Docker~Compose. 
Per un utilizzo in produzione sarebbe necessario il deployment su piattaforme cloud come AWS o Azure, sfruttando servizi gestiti quali container orchestration (ECS/EKS), database managed (MongoDB~Atlas), e broker MQTT gestiti (AWS IoT~Core). Questo garantirebbe alta disponibilità, auto-scaling in funzione del carico e ridondanza geografica.

\subsection{Miglioramento della sensibilità dell'anomaly detection}

Il modello LSTM autoencoder attuale opera sulle 15~feature grezze del dataset di telemetria refrigerata. Anomalie sottili come il \emph{sensor drift}, in cui un sensore si degrada progressivamente nel tempo senza generare un errore di ricostruzione elevato su singole osservazioni, risultano difficili da rilevare con questo approccio.

Un miglioramento consiste nell'estrazione di feature derivate da fornire in input al modello.

\subsection{Indicizzazione spaziale con PostGIS}

Il calcolo della distanza punto-polyline, utilizzato per determinare la compatibilità tra un punto di ritiro/consegna e un percorso esistente, avviene attualmente con una scansione lineare $O(n)$ su tutti i segmenti della polyline decodificata. 
Questa complessità diventa un collo di bottiglia all'aumentare del numero di viaggi e della lunghezza dei percorsi.

L'adozione di PostgreSQL con l'estensione PostGIS consentirebbe di sfruttare l'indicizzazione spaziale tramite R-tree: la funzione \texttt{ST\_Distance} opera su geometrie indicizzate con complessità $O(\log n)$, riducendo drasticamente il tempo di ricerca. Inoltre, PostGIS offre nativamente operazioni come il buffer spaziale e l'intersezione tra geometrie.

\subsection{Portale cliente per il tracking delle spedizioni}

Attualmente il ruolo \texttt{CLIENT} è limitato alla fase di ricerca e prenotazione. Una volta confermata la spedizione, il cliente non dispone di strumenti per monitorarne lo stato. Un'evoluzione prevede la creazione di un'area riservata in cui il cliente possa visualizzare lo stato della spedizione in tempo reale, con stima dell'orario di arrivo calcolata sulla base della posizione corrente del veicolo e della rotta pianificata, e notifiche automatiche all'approssimarsi del ritiro o della consegna.

\subsection{Sistema di notifiche push per il tecnico}

Il sistema attuale di notifiche per anomalie è consultabile esclusivamente tramite la dashboard web: il tecnico deve accedere alla piattaforma per verificare la presenza di nuovi alert. In un contesto operativo reale, le anomalie della catena del freddo richiedono interventi tempestivi. L'integrazione di canali di notifica push, tramite email (SMTP o servizi come SendGrid) e notifiche mobile (Firebase Cloud Messaging), garantirebbe che il tecnico venga avvisato immediatamente al verificarsi di un'anomalia, indipendentemente dal fatto che stia consultando la dashboard in quel momento.

